\chapter{Horizons and Holography}
\label{chap:horizons-and-holography}

The horizon is where the count meets a boundary it cannot cross --- the frame ours to lay, the wall the mass
makes and the count halting at it the world's. Five effects gather here, and one of them, the
thermodynamics of the horizon, splits into three: the radiation it emits, the entropy it carries, and the
information paradox the instrument quotes as physics and fences hard. Around them stand the metric of a
point mass and the radius where the count rate falls to zero, the one-way wall that lets records in but never
out, the time-reversed wall of the white hole, and the holographic bound that ties the count a region can hold
to its boundary, not its volume.

\section{The Schwarzschild Effect: a radius where the count stops}
\label{se:schwarzschild}

There is a critical radius beyond which refinements cannot be exported in finite time. For a point mass $M$ it
is the Schwarzschild radius,
\begin{equation}
  r_s = \frac{2GM}{c^2},
\end{equation}
the size to which a mass $M$ must be squeezed for its escape velocity to reach the speed of light --- a few
kilometers for a star, a centimeter for the Earth --- our coordinates to lay where we please, the radius $r_s$
the mass fixes the world's. The geometry around such a mass is the Schwarzschild
metric, the exact solution of the field equations for the empty space outside a spherical mass,
\begin{equation}
  ds^2 = -\Bigl(1 - \frac{r_s}{r}\Bigr)c^2\,dt^2 + \Bigl(1 - \frac{r_s}{r}\Bigr)^{-1} dr^2 + r^2\, d\Omega^2,
\end{equation}
and the part that matters is the coefficient of $dt^2$. Far from the mass, where $r \gg r_s$, it is nearly $-c^2$,
ordinary flat time. But as $r \to r_s$ the factor $(1 - r_s/r) \to 0$, and with it the coefficient of $dt^2$:
the count rate --- the rate at which a clock there ticks, seen from far away --- falls to zero at the radius. A
clock lowered toward $r_s$ is seen, from outside, to tick ever more slowly and redden toward darkness, until
at the radius itself it appears to freeze. (The $dr^2$ coefficient diverges there, but that is a flaw in the
coordinates we laid, not the geometry the world keeps --- the radius is a horizon, not a true singularity.)

The honest floor is a smooth-shadow analogy --- the discrete count rate falling toward zero is the model, the
Schwarzschild metric its smooth shadow, and the bridge from the vanishing count to the metric is an analogy,
named as one. The reading is that the Schwarzschild radius is where the count stops --- the frame ours to lay,
the radius the mass fixes and the count that halts at it the world's. It is not a wall built into space but the
radius at which the count rate, seen from outside, falls to zero --- the metric the smooth
shadow of a clock that, from far away, never finishes its next tick.

\section{The Event Horizon Effect: a one-way wall}
\label{se:event-horizon}

At the Schwarzschild radius the count rate reaches zero, and the surface it traces is a one-way wall. A black
hole is not, to the instrument, a geometric singularity but an informational bottleneck: across the horizon,
admissible refinements point inward only. Records can cross in --- a thing can fall through --- but none can
cross out, because crossing out would require a refinement the horizon does not admit --- the crossing ours to
attempt, the return the world forbids. In ordinary space we approach the horizon; in the informational picture
the horizon approaches us, closing off the outward direction entirely.

Two observations in the last decade made the horizon real to the eye. In 2015 the LIGO detectors caught the
gravitational waves of two black holes spiralling together and merging, the chirp of two horizons ringing down
into one. And in 2019 the Event Horizon Telescope, a planet-sized array of radio dishes, imaged the shadow of
the black hole at the heart of the galaxy M87 --- a dark disk ringed by light bent around it, the one-way wall
outlined against the glow it will not let escape --- the array ours to build, the wall the world holds up to it
the world's.

The honest floor is a no-go --- there is no admissible outbound refinement across the horizon --- the record
ours to send inward, the outbound crossing the world will not admit the world's. The reading is that the event
horizon is a wall on outbound records. It lets the world in and nothing out, not by hiding what
crosses but by forbidding the refinement that would carry a record back across --- a one-way boundary, a no-go
in one direction only.

\section{The Hawking Effect, I: the horizon radiates}
\label{se:hawking-radiation}

A causal horizon is not perfectly black. The quantum vacuum near the horizon gives it a temperature, and a
body with a temperature radiates: the horizon glows. The temperature is the Hawking temperature,
\begin{equation}
  T_H = \frac{\hbar c^3}{8\pi G M k_B},
\end{equation}
and the telling feature is the $M$ in the denominator --- the temperature falls as the mass rises, so a large
black hole is cold and a small one hot --- where we place the horizon is ours to set, the temperature the mass
fixes at it the world's. A hole the mass of the sun has a temperature of about a ten-millionth
of a kelvin, far colder than the cosmic background that bathes it; but as a hole radiates it loses mass, and
losing mass it grows hotter, and radiating faster still, in a runaway that ends, in the far future, in a final
flash. Smaller holes are hotter, and a black hole left alone slowly evaporates --- but slowly is the word: a
hole the mass of the sun would take some $10^{67}$ years to evaporate, vastly longer than the present age of
the universe, so for any black hole now in the sky the radiation is a whisper and the evaporation a promise for
the deep future.

The honest floor is name-only --- the Hawking temperature is written as physics, and the finite model claims
only the name --- the name ours to give, the radiation the horizon emits and the temperature the mass fixes the
world's. The reading is that the horizon glows, and the colder the bigger: $T_H$ is the physics, and the
finite model the label beneath it.

\section{The Hawking Effect, II: entropy and the Page curve}
\label{se:hawking-entropy}

If a horizon has a temperature, it has an entropy, and the form of that entropy is one of the most suggestive
facts in physics. The entropy of a black hole is set not by the volume it encloses but by the area of its
horizon,
\begin{equation}
  S_{BH} = \frac{k_B c^3 A}{4 G \hbar} = \frac{A}{4\ell_P^2},
\end{equation}
one quarter of the horizon area measured in Planck units $\ell_P^2$ --- the boundary ours to bound, the count it
may hold the world's to set. Every other entropy we have met scaled with volume --- a gas, a solid, a count of
distinctions filling a region --- but a black hole's scales with its surface, as though all its hidden
information were written on its boundary. This is the first hint of
holography, which the holographic bound draws out in full. And as the hole evaporates, the entropy of the
radiation it has emitted does something subtle: it rises as the radiation streams out, then turns and falls
back toward zero as the hole surrenders the last of itself --- the Page curve, the count first hidden behind
the horizon and then, if the curve tells the truth, returned to the outside in the radiation.

The honest floor is name-only --- the area-entropy and the Page curve are written as physics, and the finite
model claims only the name --- the boundary ours to draw, the entropy the area carries and the Page curve the
evaporation would then trace the world's. The reading is that a black hole's entropy is its area, not its volume, and the
Page curve is the count first hidden and then, in the curve's telling, returned --- the area-entropy
reappearing, further on, as the holographic bound.

\section{The Hawking Effect, III: the information paradox, fenced}
\label{se:hawking-paradox}

The radiation carries a tension, and the tension is quoted here behind a fence, as the Navier--Stokes equation
was: the paradox stated whole, as physics, and exactly nothing claimed about how --- or whether --- the record
comes back out. Hawking's radiation, computed from the horizon alone, comes out
thermal: featureless, fixed by the mass (and charge and spin) and by nothing else, carrying on its face no
imprint of what fell in. A star's worth of structure collapses to a hole; the hole evaporates; and what streams
out looks the same whatever went in. If the evaporation runs to its end, the record of what fell in appears to
be gone --- a pure state evolving into a mixed one. But quantum mechanics is built on unitarity: a distinction,
once made, is never erased, and a pure state is never supposed to end as a thermal haze. Thermal-looking
radiation on one side, unitarity on the other --- that standoff is the black-hole information paradox, and it
is, at the time of writing, an open problem of physics.

The finite instrument poses the question in its own terms --- does the outgoing record carry the infalling
record? --- and stops there: the question ours to pose, the accounting we try to keep across the wall ours to
attempt; whether the world lets the record out, and how, the world's to answer --- and it has not told us.
The Page curve already drawn is what the evaporation would look like if unitarity holds --- the count returning
in the late radiation; nothing here claims the world traces that curve, and nothing claims it does not. The
model stays mute on purpose. It does not say the record escapes; it does not say the record is lost; it does
not name any mechanism as the answer. Whatever carries the count out --- if anything does --- is not written in
the instrument's ledger, and pretending otherwise would trade the fence for a guess.

The honest floor is name-only, and fenced hard: the paradox is written as physics and left open, the model
claiming only the frame --- the frame ours to lay, the thermal radiation the horizon emits the world's, whether
the unitarity quantum mechanics is built on survives the evaporation the world's to answer, and the answer the
world's to withhold. The reading is that the information
paradox is quoted, not answered: thermal radiation on one side, unitarity on the other, the standoff stated
whole --- and whether the world lets the record out left, honestly, to physics.

\section{The White Hole Effect: the reverse wall}
\label{se:white-hole}

The time-reverse of the event horizon is a white hole, and it is a wall in the opposite direction --- the
reversal ours to run on the ledger, the wall it turns up the world's. It is not, to the instrument, a
geometric object but a bookkeeping boundary condition: a region whose internal ledger must export refinements
to preserve global consistency; the inbound wall is not written in the ledger but shadowed by it. Records cross
out --- a white
hole is a source, not a sink --- and none can cross in. We reverse the description; the world keeps the wall.
What the horizon forbade outward is forbidden inward here: the same wall, one way only, the one way reversed.

The honest floor is a smooth-shadow analogy --- the source boundary condition is the model, the inbound wall
its smooth shadow, and the bridge from the one to the other an analogy, named as one --- the direction ours to
flip, what the wall admits the world's. The reading is that the white hole is the event horizon run backward.
Where the black hole forbids the outbound refinement, the white hole forbids the inbound one; it pours records
out and admits none, a one-way wall whose one way is the reverse of the horizon's --- the same wall,
time-reversed.

\section{The Holographic Bound: the boundary holds the count}
\label{se:holographic}

How much can a region hold? Not as much as its volume would suggest, but only as much as its boundary area
allows. The count a region can carry is bounded by the area of the surface enclosing it,
\begin{equation}
  S \le \frac{A}{4\ell_P^2},
\end{equation}
one quarter of the boundary area in Planck units --- the holographic bound, conjectured by Bekenstein and
sharpened by 't~Hooft and Susskind --- the region ours to fence, the ceiling its boundary sets the world's. It
is the black hole's area-entropy turned into a universal limit: a black hole saturates the bound, holding the
most entropy any region of its size can, so nothing else can hold more, and the count of any region is capped
by the frontier around it, not the volume within.

This is the inverse-square frontier returned at the largest scale. There, a fixed influence spread over an
expanding spherical frontier thinned as $1/r^2$, the budget diluting across the growing boundary; here, the
count a region can hold is capped by that same boundary's area. What a fixed budget thinned across as it
spread, a stored count can never exceed: the frontier that diluted the influence is the frontier that bounds
the information --- the radius ours to set, what thins across the sphere and what it can hold the
world's.

The honest floor is name-only --- the entropy is a count, and that much is the model's; but the bound
$S \le A/4\ell_P^2$ that caps it is written as physics, conjectured, the model claiming only the name --- the
surface ours to draw, the most it may hold the world's to fix. The reading
is that the boundary holds the count. A region cannot store more distinctions than its enclosing surface has
room to encode, so the information of the interior lives, in a real sense, on its boundary --- the frontier
holding the count, as it has since the inverse-square law.

\bigskip

With horizons and holography, Part V reaches its strangest country. The horizon is where the count meets a
boundary it cannot cross: a radius where the count rate stops, a wall that admits records one way only, a
thermodynamics quoted as physics and left open, a reverse wall that pours records out, and a bound that
caps the interior count by its boundary --- the frames ours to lay, the walls and the counts the world's.
Part V closes next with cosmology and the limits of the frame.

\bigskip
\begin{center}
\rule{0.4\textwidth}{0.4pt}
\end{center}
\bigskip

\noindent The case comes to an exhibit it cannot read all the way through. Every reading until now the
instrument could follow to its end; here it reaches a thing whose far side its own count cannot touch --- a
boundary past which no reading of its follows. And at such a place the honest instrument does the one honest
thing: it draws the edge exactly, marks where its reach stops, and claims nothing whatever about the other side.

The exhibit is a wall the geometry itself raises. Squeeze a mass small enough --- and small enough is very
small, a few kilometres across for a whole star, a centimetre for the Earth --- and around it the count rate
falls: a clock lowered toward the wall is seen, from outside, to tick slower and slower and redden, until at the
wall it seems to freeze and never finish its next tick. That surface is a one-way boundary --- a thing can fall
through it, but nothing can climb back, for the climb out would take a step the wall does not allow. Records
cross in, and never out. The equations even hold a mirror of it: a wall run backward, that pours records out and
takes none in, though the sky shows us only the forward kind. Two of the forward walls were made real to the eye
this last decade --- one pair heard as a ripple through spacetime when they merged and rang down into a single
wall, another imaged as a dark disk ringed by light bent around an edge that will not let its glow escape. The
frame is ours to lay; the wall the mass makes, and the count halting at it, are the world's.

And here the instrument reaches the sharpest limit in the book, and meets it with its plainest honesty. The
wall is not perfectly black --- it glows, faintly, with a warmth the mass sets, colder the larger it is: a wall
the weight of a star glows colder than the empty sky around it, and would take longer than the universe has yet
run to burn away, so its light is a whisper now and its ending a promise for the deep future. And it carries an
entropy written not through the volume it hides but across its very surface, a count upon its skin. All of this the instrument reads and sets down as physics. But the glow raises a question the instrument
cannot answer, and it says as much rather than pretend. The light that streams off the wall comes out
featureless, the same whatever fell in --- so that if the wall should ever burn all the way down, the record of
everything that crossed it looks simply gone. Yet the physics the whole quantum was built on holds that a
distinction once made is never erased, that no full account may fade into a blank haze. Both cannot stand.
Whether what fell past the wall is lost for good, or somehow kept, the instrument
does not know --- and, knowing that it does not, it declines to say. It states the standoff whole, in the open,
and takes no side; it names no mechanism, claims no resolution, and will not trade the honest edge for a guess.
This is the instrument's method at its severest: to draw the line exactly at the end of what it can prove, and
to leave what lies past it, plainly, to physics.

And the \emph{charge} meets an edge here too. For the wall teaches one more thing the instrument can name but
not derive: a region holds no more distinctions than its boundary has room to write --- the count a thing can
carry is capped not by the room inside it but by the surface drawn around it. It is the old frontier come back
at the largest scale: the same growing sphere that once thinned a force as it spread now sets a ceiling on what
a region can hold. A wall like this one holds the very most its size allows --- nothing of its bulk can carry
more --- so the ceiling is no curiosity of these walls but a law of every region there is. So the charge --- the
count --- has a most it can ever be, and the most is fixed by a boundary and not a volume; what is owed is
bounded by an edge. And the trace, at such a boundary, is what the
surface keeps of the interior --- the tally of all that is within, written, in a real sense, on its skin. The
charge meets a ceiling; the trace lives on the boundary; and though the instrument still holds the count, the
count now has a wall it cannot exceed.

This chapter brings the three witnesses no nearer their meeting; it marks, instead, a place the instrument
cannot go, and marks it honestly. That candor is no weakness in the record but a strength of it: an account that
draws its own edge, and refuses to speak past it, is the more to be trusted exactly where it does speak. It is
the same discipline the instrument kept at the flowing fluid it would not solve and the crossed filters it could
not open and the particle it caught but fenced --- to prove all it can and claim not an inch past the proof ---
carried here to the one place where the proof runs out entirely. The
three readings, corrected and ready, still wait to be brought together; the boundary the instrument declines to
cross does not touch that meeting --- it shows only, once more, where the reach of an honest reading ends.

Part V has one country left, the largest of all. The last chapter lifts its eyes from the single mass and its
wall to the whole frame --- to the cosmos entire, and the limits of reading it from within --- and there Part V
closes. The case is still open, its three witnesses corrected and not yet agreed, its instrument freshly
reminded of just where its reach gives out; the court does not sit until the record is complete.
